% !TeX root = ../教材题库_全部条目_转写工作稿.tex
% Source: XJ-B1-C06-CA-P226-Y0277
% 第6章 / 6.3 统计图表 / 教材 p.226 / PDF 第233页
% OCR未校对：公式与文字均可能有误；未完成数学TeX转写；禁止直接发给学生
% 原题及相邻原解截图（核对用）：../题目原图/XJ-B1-C06-CA-P226-Y0277-01.webp
% 本文件只在转写完成后作为题干使用；OCR 条目尚不是可靠题干。
\begin{bankitem}{6.3 案例 案例2}{第6章 / 6.3 统计图表 / 教材 p.226 / PDF 第233页}{XJ-B1-C06-CA-P226-Y0277}{OCR未校对：公式与文字均可能有误；未完成数学TeX转写；禁止直接发给学生}
\begin{ocrdraft}
[教材p.226 / PDF第233页]
\par
案例2甲、乙两家企业近八年对研究与试验发展的资金投人如表6一2所示，
\par
试根据数据绘制折线统计图，并分析两家企业资金投人情况．
\par
表6一2
\par
时
\par
间
\par
甲企业资金投人的数量/万元
\par
乙企业资金投人的数量/万元
\par
第一年
\par
31.5
\par
105
\par
第二年
\par
35
\par
120
\par
第三年
\par
39．8
\par
134．5
\par
第四年
\par
45．4
\par
150
\par
第五年
\par
4\&3
\par
16\&7
\par
第六年
\par
53
\par
18\&9
\par
第七年
\par
54．7
\par
218
\par
第八年
\par
57
\par
243．2
\par
解用横轴表示时间，纵轴表示资金投人的数量，根据表中数据描点，然后分
\par
别用线段依次连接表示甲、乙企业的点，就得到折线统计图，如图6，3一3所示．
\par
资金投人的数量/万元
\par
300
\par
\_甲企业
\par
250
\par
218
\par
一乙企业
\par
200
\par
158．7
\par
150
\par
150
\par
120
\par
105
\par
100
\par
54．7
\par
50
\par
45．4
\par
0
\par
第一年第二年第三年第四年第五年第六年第七年第八年时间
\par
图6．3一3
\par
从图6．3一3可以清楚地看出，甲、乙两家企业对研究与试验发展的资金投人数
\par
量逐年增加；近八年，乙企业对研究与试验发展的资金投人数量均超出甲企业，并
\par
且这种差距有扩大的趋势．
\end{ocrdraft}
\end{bankitem}
